\section{Introduction}
\label{sec:introduction}

Thompson's group \(F\) consists of the increasing piecewise linear
homeomorphisms of \([0,1]\) with finitely many pieces, dyadic rational
breakpoints, and slopes in \(2^{\Z}\). We use the multiplication convention
\(hg=h\circ g\) and regard \(F\) as a discrete group.

A discrete group \(G\) is \emph{amenable} if its bounded real-valued
functions admit a positive normalized left-invariant mean: a linear
functional \(M:\ell^\infty(G;\R)\to\R\) satisfying \(M(1)=1\),
\(M(\varphi)\ge0\) whenever \(\varphi\ge0\), and
\(M(g\mapsto\varphi(hg))=M(\varphi)\) for every \(h\in G\).
This invariant-mean viewpoint belongs to the classical theory developed
by von Neumann and Day \cite{vonNeumann1929,Day1957}.
We use the following forward implication of the F\o lner criterion
\cite{Folner1955}: if \(G\) is amenable, then for every finite
\(S\subset G\) and every \(\varepsilon>0\), there is a nonempty finite
\(A\subset G\) such that
\begin{equation}\label{eq:folner}
  \frac{|hA\mathbin{\triangle}A|}{|A|}<\varepsilon
  \qquad(h\in S).
\end{equation}
The modern left-translation formulation in
\cite[Theorem~1.2]{Moore2013} bounds the sum of these ratios and hence
implies each strict individual bound by taking a smaller tolerance.
Thus amenability requires finite sets
whose relative boundaries are simultaneously small under any fixed
finite family of translations.

Thompson introduced \(F\) in 1965; an early published construction appears
in McKenzie and Thompson \cite{McKenzieThompson1973}, with the identification
explained by Cannon, Floyd, and Parry \cite[pp.~215--216]{CannonFloydParry1996}.
The latter authors record Geoghegan's 1979 conjecture that \(F\) is
nonamenable \cite[p.~227]{CannonFloydParry1996}. We prove this conjecture.

\begin{theorem}\label{thm:main}
Thompson's group \(F\) is not amenable.
\end{theorem}

\paragraph{Context and prior work.}
Two classical structural facts help explain the difficulty of the
amenability problem. Brin and Squier show that \(F\) contains no
nonabelian free subgroup \cite[Theorem~3.1]{BrinSquier1985}, so the familiar
free-subgroup obstruction does not apply. Yet \(F\) is not elementary
amenable \cite[Theorem~4.10]{CannonFloydParry1996}: it lies outside the class
generated by finite and abelian groups under subgroups, quotients,
extensions, and directed unions. Neither fact decides ordinary
amenability. The group also has strong topological finiteness properties:
Brown and Geoghegan construct a classifying space with finitely many
cells in each dimension \cite{BrownGeoghegan1984}.

Finite approximations have provided several ways to study the remaining
question. Moore proves tower lower bounds on the sizes of possible
F\o lner sets \cite[Theorem~1.1]{MooreFolner2013} and characterizes
amenability of \(F\) by a scalar convex Ramsey property for finite rooted
ordered binary trees \cite[Theorem~3.1]{Moore2013}. Guba improves densities
of finite Cayley subgraphs \cite{Guba2025Density} and derives restrictions
on right-invariant means from a partition into seven diagram-defined
classes \cite{Guba2026}. From an operator-algebraic
direction, Haagerup and Olesen prove that simplicity of the reduced group
\(C^*\)-algebra of Thompson's group \(T\) would imply nonamenability of
\(F\) \cite[Theorem~5.5]{HaagerupOlesen2017}. These results expose constraints
on amenability through finite sets, trees, diagrams, and representations.
The proof below uses finite dyadic partitions and scalar correlations of
Hilbert-valued colors.

There have also been claims of a complete resolution in both directions.
Akhmedov's 2021 version claims nonamenability through a height-function
criterion \cite{Akhmedov2021}. Shavgulidze claimed amenability
\cite{Shavgulidze2009}; Moore identified errors in that approach
\cite{Moore2011Critique}. Moore separately withdrew an amenability claim
after Akhmedov identified an error in Lemma~4.13 of that manuscript
\cite{Moore2012Withdrawn}. That withdrawal concerns a different work from
his published Ramsey characterization cited above.

\paragraph{Proof strategy.}
The proof turns approximate translation invariance of finite averages of
scalar functions on \(F\) into an approximate fixed point of a Lipschitz
map on a Hilbert ball. The
following infinite-dimensional phenomenon supplies a map for which that
conclusion is impossible.

\begin{lemma}[Benyamini--Sternfeld]\label{lem:displacement}
There are a real Hilbert space \(H\), a Lipschitz map \(f:B\to B\) on its
closed unit ball, and a constant \(\delta>0\) such that
\[
  \norm{f(x)-x}\ge\delta\qquad(x\in B).
\]
\end{lemma}

Benyamini and Sternfeld prove this assertion for every infinite-dimensional
normed space \cite[Theorem, part~(3), p.~439]{BenyaminiSternfeld1983}.
For completeness, Appendix~\ref{app:auxiliary} constructs such a map on
\(L^2([0,1];\R^2)\) with \(\delta=1/2\), including all the needed
Lipschitz estimates.

Fix this map \(f\), a positive Lipschitz constant \(L\), and a large
integer \(D\). A basic dyadic interval is a cell of a uniform dyadic
partition, and a basic partition is a finite partition into such cells.
Choose \(D\) basic intervals with positive gaps between them and endpoints
in \((0,1)\), called parents. Put an affine copy of this family inside
each parent, calling the copies descendants. When a parent is a union
of partition cells, restricting to it and stretching back to \([0,1]\)
gives its normalized restriction. Color each basic partition recursively
by applying \(f\) to the mean of its parent-restriction colors when all
these restrictions are defined, and by zero otherwise. Proper
restriction decreases the number of cells, so the recursion is well
founded. For a sufficiently fine image partition under \(g\in F\), its
parent colors \(X_i\) and descendant colors \(Y_{ij}\) therefore satisfy
\[
 m=\frac1D\sum_{i=1}^D X_i,\qquad
 z_i=\frac1D\sum_{j=1}^D Y_{ij},\qquad
 X_i=f(z_i)\quad(1\le i\le D).
\]

The group \(F\) can carry any ordered pair of separated basic intervals
with endpoints in \((0,1)\) affinely onto any other such pair. Choose a
finite set \(S\) of these transports comparing every separated pair among
the parents and descendants with one reference pair. The set \(S\) is
fixed before the arbitrary nonempty
finite set \(A\subset F\); only the common partition level, chosen fine
enough for \(A\cup SA\), may depend on \(A\). Write \(\eta\) for the
largest relative \(S\)-boundary of \(A\). Exact covariance of normalized
restrictions and finite-average cancellation make all averaged scalar
correlations of separated interval colors differ from one common value
by at most \(\eta\).

Expand the average over \(A\) of \(\norm{z_i-m}^2\). The common
correlation cancels; diagonal and nested pairs occupy only a fraction
\(1/D\) of each relevant sum, as Figure~\ref{fig:mixed-pairs} illustrates.
Unit-ball bounds control those exceptions, giving a universal constant
times \(D^{-1}+\eta\). Since \(m\) is also the mean of the \(f(z_i)\),
convexity and Lipschitz continuity bound the average squared displacement
\(\norm{m-f(m)}^2\) by a constant times \(L^2(D^{-1}+\eta)\).
For large \(D\) and small \(\eta\), this contradicts the uniform lower
bound \(\delta^2\). Section~\ref{sec:proof} makes the argument precise and proves a
uniform positive boundary bound in Proposition~\ref{avg:boundary},
contradicting \eqref{eq:folner}.

Section~\ref{sec:consequences} records consequences for uniformly bounded
representations and percolation, using separate companion theorems.
